\begin{tikzpicture}[
    thick, font=\small,
    contact/.style={line width=1pt},
    coil/.style={circle, draw, line width=1pt, minimum size=6mm, inner sep=0pt},
    block/.style={draw, rectangle, minimum width=2.4cm, minimum height=1cm,
                  align=center}
]

% --- Szyny zasilające ---
\draw (0,0) -- (0,10);
\draw (15,0) -- (15,10);
\node[above] at (0,10) {$L$};
\node[above] at (15,10) {$N$};

% =========== Szczebel 1: P = S1 || S2 || S3 || S4 ===========
\draw (0,9) -- (1,9);
\draw (1,9) -- (1,5.5);
\foreach \y/\lbl in {9/S_1, 8/S_2, 7/S_3, 6/S_4} {
    \draw (1,\y) -- (1.6,\y);
    \draw[contact] (1.6,\y-0.2) -- (1.6,\y+0.2);
    \draw[contact] (2.1,\y-0.2) -- (2.1,\y+0.2);
    \draw (2.1,\y) -- (3,\y);
    \node[above] at (1.85,\y+0.2) {$\lbl$};
}
\draw (3,9) -- (3,6);
\draw (3,7.5) -- (5,7.5);
\node[above] at (5,7.5) {$P$};
\draw (5,7.5) -- (7,7.5);

% =========== Szczebel 2: przerzutnik T z kasowaniem ===========
\draw (0,4) -- (1,4);
\draw (1,4) -- (1,2);

% Gałąź 1: Q * /P
\draw (1,4) -- (1.6,4);
\draw[contact] (1.6,3.8) -- (1.6,4.2);
\draw[contact] (2.1,3.8) -- (2.1,4.2);
\node[above] at (1.85,4.2) {$Q$};
\draw (2.1,4) -- (2.8,4);
\draw[contact] (2.8,3.8) -- (2.8,4.2);
\draw[contact] (3.3,3.8) -- (3.3,4.2);
\draw (2.8,3.8) -- (3.3,4.2);
\node[above] at (3.05,4.2) {$\bar P$};
\draw (3.3,4) -- (4.5,4);

% Gałąź 2: /Q * P
\draw (1,2) -- (1.6,2);
\draw[contact] (1.6,1.8) -- (1.6,2.2);
\draw[contact] (2.1,1.8) -- (2.1,2.2);
\draw (1.6,1.8) -- (2.1,2.2);
\node[above] at (1.85,2.2) {$\bar Q$};
\draw (2.1,2) -- (2.8,2);
\draw[contact] (2.8,1.8) -- (2.8,2.2);
\draw[contact] (3.3,1.8) -- (3.3,2.2);
\node[above] at (3.05,2.2) {$P$};
\draw (3.3,2) -- (4.5,2);

% Połączenie gałęzi
\draw (4.5,4) -- (4.5,2);
\draw (4.5,3) -- (6,3);

% Styk /T
\draw[contact] (6,2.8) -- (6,3.2);
\draw[contact] (6.5,2.8) -- (6.5,3.2);
\draw (6,2.8) -- (6.5,3.2);
\node[above] at (6.25,3.2) {$\bar T$};
\draw (6.5,3) -- (8,3);

% Cewka Q
\node[coil] (Q) at (8.6,3) {$Q$};
\draw (8,3) -- (Q.west);
\draw (Q.east) -- (15,3);

% =========== Blok TON (timer 2 min) ===========
\draw (0,1) -- (1,1);
\draw (1,1) -- (7,1);
\node[block] (TON) at (8.5,1) {TON (\SI{2}{\minute})};   % <-- TU bez \\
\draw (7,1) -- (TON.west);
\draw (TON.east) -- (15,1);
\node[above] at (4,1) {$M$ (czujnik ruchu)};

% Sygnał T z timera
\node at (11,1.6) {$T$};
\draw[dashed,->] (9.5,1.5) -- (6.25,2.7);

\end{tikzpicture}